% =============================================================================
%  Atto 5 — The Upstream Path & the Contract  (~10 min, ~6 slide)  [Parte II.B]
%  The same path, walked backwards: from encoders to /odom, then DBC.
% =============================================================================

\section{Part II.B — Upstream: closing the loop}

% ---------------------------------------------------------------------------- %
\begin{frame}{The same chain, walked backwards}
  \centering
  \resizebox{0.96\textwidth}{!}{%
  \begin{tikzpicture}[
    node distance=4mm,
    upstreamstep/.style={minimum width=28mm, minimum height=10mm, align=center},
    backarr/.style={line width=0.65pt,
      -{Latex[length=1.8mm,width=1.4mm]}, shorten <=0.7mm, shorten >=0.7mm}
  ]
    \node[ll, upstreamstep]                (mcu) {STM32\\\scriptsize encoders};
    \node[bus, upstreamstep, right=of mcu] (bus) {socketcan\\\scriptsize CAN frame};
    \node[hl, upstreamstep, right=of bus]  (rpt) {\texttt{report\_parser}\\\scriptsize decode};
    \node[hl, upstreamstep, right=of rpt]  (sys) {\texttt{MiviaRoverSystem}\\\scriptsize watchdog};
    \node[hl, upstreamstep, right=of sys]  (dd)  {\texttt{diff\_drive}\\\scriptsize FK + odom};
    \node[hl, upstreamstep, right=of dd]   (odom){\texttt{/odom}};

    \draw[arrll, backarr] (mcu.north) -- ++(0,4mm)
      -- node[above,msg]{RPMs} ($(bus.north)+(0,4mm)$) -- (bus.north);
    \draw[arrhl, backarr] (bus.south) -- ++(0,-4mm)
      -- node[below,msg]{\texttt{can\_rx}} ($(rpt.south)+(0,-4mm)$) -- (rpt.south);
    \draw[arrhl, backarr] (rpt.north) -- ++(0,4mm)
      -- node[above,msg]{\texttt{EncoderRpms}} ($(sys.north)+(0,4mm)$) -- (sys.north);
    \draw[arrhl, backarr] (sys.south) -- ++(0,-4mm)
      -- node[below,msg]{state if.} ($(dd.south)+(0,-4mm)$) -- (dd.south);
    \draw[arrhl, backarr] (dd.north) -- ++(0,4mm)
      -- node[above,msg]{\texttt{Odometry}} ($(odom.north)+(0,4mm)$) -- (odom.north);
  \end{tikzpicture}}

  \vspace{0.6em}
  \begin{itemize}\setlength\itemsep{1pt}
    \item Same rings as Atto~4, now turning encoder feedback into \texttt{/odom}.
    \item The only places where meaning changes: \emph{decode},
          \emph{watchdog}, \emph{odometry}.
    \item The DBC is the byte-level contract underneath both directions.
  \end{itemize}
\end{frame}

\begin{comment}
% ---------------------------------------------------------------------------- %
\begin{frame}[fragile]{Ring 1 — the firmware emits RPMs}
  \begin{columns}[T,onlytextwidth]
    \column{0.55\textwidth}
      Every \(5\)~ms the firmware emits one \texttt{EncoderRpms} frame.

      \vspace{0.3em}
      The loop is just:
      \begin{enumerate}\setlength\itemsep{2pt}
        \item sample the four QEI counters,
        \item compensate roll-over and convert ticks \(\to\) RPM,
        \item pack four signed \(16\)-bit values and enqueue CAN.
      \end{enumerate}

      \vspace{0.4em}
      \footnotesize
      \textbf{Frame:} ID \texttt{0x0B}, \(8\)~B, scale \(1/182\) RPM.\\
      \textcolor{black!60}{This is the deterministic end of the pipeline:
      sampling, conversion, packing, done.}

    \column{0.45\textwidth}
      \begin{lstlisting}[language=, basicstyle=\ttfamily\scriptsize]
# DBC excerpt (rover.dbc)
BO_ 11 EncoderRpms: 8 Rover
  SG_ FrontLeft  :  0|16@1- (1,0) "rpm"
  SG_ RearLeft   : 16|16@1- (1,0) "rpm"
  SG_ FrontRight : 32|16@1- (1,0) "rpm"
  SG_ RearRight  : 48|16@1- (1,0) "rpm"
      \end{lstlisting}

      \scriptsize\textcolor{black!60}{One frame, four fields, one scale.
      That is the upstream protocol.}
  \end{columns}
\end{frame}
\end{comment}

% ---------------------------------------------------------------------------- %
\begin{frame}[fragile]{Ring 1 — \texttt{report\_parser}: from bytes to types}
  \begin{columns}[T,onlytextwidth]
    \column{0.45\textwidth}
      \texttt{report\_parser} subscribes to \texttt{input/can\_rx}
      and dispatches by CAN ID:

      \begin{itemize}\setlength\itemsep{2pt}
        \item \texttt{0x0B} $\to$ \texttt{EncoderRpms}\\
              {\footnotesize\textcolor{black!60}{scale: \(1/182\)}}
        \item \texttt{0x0C} $\to$ \texttt{ImuGyroXy}
        \item \texttt{0x0D} $\to$ \texttt{ImuAccelXy}
        \item \texttt{0x0E} $\to$ \texttt{ImuZ}
      \end{itemize}

      \textbf{Stale-data guard.}\;
      At \(50\)~Hz, publish only fresh decoded frames
      (\texttt{report\_timeout\_ms}); otherwise log a throttled error.

    \column{0.50\textwidth}
      \begin{lstlisting}[language=C++, basicstyle=\ttfamily\scriptsize]
// Pseudocode -- per-ID dispatch
switch (frame.id) {
  case EncoderRpms::ID:
    last_rx_enc = now();
    decode(frame.data, &enc_entity);
    enc_msg.fl = enc_entity.fl * (1.0/182.0);
    enc_msg.rl = enc_entity.rl * (1.0/182.0);
    enc_msg.fr = enc_entity.fr * (1.0/182.0);
    enc_msg.rr = enc_entity.rr * (1.0/182.0);
    break;

  case ImuGyroXy::ID:   /* ... */ break;
  case ImuAccelXy::ID:  /* ... */ break;
  case ImuZ::ID:        /* ... */ break;
}
      \end{lstlisting}
      \scriptsize\textcolor{black!60}{Decoding/scaling are mechanical: all the
      ``intelligence'' is in the DBC. Adding a new feedback message =
      adding a case here.}
  \end{columns}
\end{frame}

% ---------------------------------------------------------------------------- %
\begin{frame}[fragile]{Ring 2 — \texttt{read()}: RT pull \(+\) watchdog}
  \begin{columns}[T,onlytextwidth]
    \column{0.55\textwidth}
      Each \texttt{read()} tick pulls the latest encoder sample from the
      RT buffer:

      \begin{enumerate}\setlength\itemsep{2pt}
        \item Check age: \(t_{\text{now}} - t_{\text{enc}}\).
        \item Late frame \(>150\)~ms: zero velocity, count timeout.
        \item Too many timeouts: \texttt{fault\_stop\_}, return
              \texttt{ERROR}.
        \item Fresh frame: RPM $\to$ rad/s, update velocity + position.
      \end{enumerate}

      \vspace{0.2em}
      \alert{Why graceful degradation?}\;
      One late frame slows us down; persistent silence stops the run.

    \column{0.45\textwidth}
      \begin{lstlisting}[language=C++, basicstyle=\ttfamily\scriptsize]
// Pseudocode -- read()
const EncoderSample* enc =
    encoder_buffer_.readFromRT();
if (!enc || !enc->valid) return OK;

age = (time - enc->stamp);
if (age > feedback_timeout_) {
  consecutive_timeouts_++;
  for (k=0..3) hw_velocities_[k] = 0;
  if (consecutive_timeouts_ >= max_to)
    return ERROR;       // hard fault
  fault_stop_ = true;   // soft fault
  return OK;
}

consecutive_timeouts_ = 0;
fault_stop_           = false;
for (k=0..3) {
  w = rpm_to_rad_s(enc->rpm[k]);
  hw_velocities_[k]  = w;
  hw_positions_[k]  += w * dt;
}
      \end{lstlisting}
  \end{columns}
\end{frame}

% ---------------------------------------------------------------------------- %
\begin{frame}{Ring 3 — \texttt{diff\_drive\_controller} produces \texttt{/odom}}
  \begin{columns}[T,onlytextwidth]
    \column{0.50\textwidth}
      Once the four wheel velocities are in the state interfaces, the
      \texttt{diff\_drive\_controller} closes the loop \emph{for free}:

      \begin{itemize}\setlength\itemsep{2pt}
        \item Aggregates the two left and two right wheels into a 2-side
              differential model.
        \item Applies the inverse of the kinematics it used downstream:
              \[
                v_x = \tfrac{R}{2}(\omega_L+\omega_R),\quad
                \omega_z = \tfrac{R}{L}(\omega_R-\omega_L)
              \]
        \item Integrates \((v_x,\omega_z)\) into a planar pose; publishes
              \texttt{nav\_msgs/Odometry} on \texttt{/odom} \(+\) the
              \texttt{odom $\to$ base\_link} TF.
      \end{itemize}

    \column{0.45\textwidth}
      \textbf{Frames \(+\) configuration.}
      \begin{itemize}\setlength\itemsep{1pt}
        \item \texttt{odom\_frame\_id: odom}
        \item \texttt{base\_frame\_id: base\_footprint}
        \item \texttt{enable\_odom\_tf: true}
        \item \texttt{publish\_rate: 50.0}~Hz
      \end{itemize}
  \end{columns}
\end{frame}

% ---------------------------------------------------------------------------- %
\section{The contract: a single DBC file}

\begin{frame}{The DBC: one file, two stacks}
  \centering
  \begin{tikzpicture}[node distance=8mm]
    \node[bus, text width=24mm, minimum height=18mm, align=center] (dbc)
      {\textbf{rover.dbc}\\\scriptsize 5 messages\\\scriptsize 1\,Mbit/s};
    \node[hl, text width=45mm, minimum height=18mm, align=center, left=of dbc] (sw)
      {ROS~2 stack\\
      
      \scriptsize \texttt{MiviaRoverSystem}\\
      \scriptsize \texttt{control\_command}\\
      \scriptsize \texttt{report\_parser}};
    \node[ll, text width=45mm, minimum height=18mm, align=center, right=of dbc] (fw)
      {Firmware stack\\\scriptsize ISR + \texttt{can\_manager}\\
       \scriptsize \texttt{can\_msg\_pack}\\
       \scriptsize \texttt{unpack\_*}};

    \draw[arr,<->,line width=0.7pt,shorten <=1mm,shorten >=1mm] (sw.east)  -- (dbc.west);
    \draw[arr,<->,line width=0.7pt,shorten <=1mm,shorten >=1mm] (dbc.east) -- (fw.west);
  \end{tikzpicture}

  \vspace{0.4em}
  {\scriptsize
  \begin{tabular}{@{}lcll@{}}
    \toprule
    Message       & ID    & Direction        & Fields\\
    \midrule
    \texttt{Reference}    & \texttt{0x0A} & Jetson $\to$ Rover & 4$\times$ \texttt{int16} RPM\\
    \texttt{EncoderRpms}  & \texttt{0x0B} & Rover  $\to$ Jetson& 4$\times$ \texttt{int16} RPM ($\div$182)\\
    \texttt{ImuGyroXy}    & \texttt{0x0C} & Rover  $\to$ Jetson& 2$\times$ \texttt{int32} \([\text{rad/s}]\)\\
    \texttt{ImuAccelXy}   & \texttt{0x0D} & Rover  $\to$ Jetson& 2$\times$ \texttt{int32} \([\text{m/s}^2]\)\\
    \texttt{ImuZ}         & \texttt{0x0E} & Rover  $\to$ Jetson& gyro\(_z\) + accel\(_z\)\\
    \bottomrule
  \end{tabular}}

  \vspace{0.3em}
  \footnotesize
  \alert{The \texttt{.dbc} is the only surface of agreement} between SW and
  FW. Anything else evolves independently. Human-readable, git-diffable,
  parseable by every CAN tool out there.
\end{frame}

% ---------------------------------------------------------------------------- %
\begin{comment}
\begin{frame}{Why this pattern matters beyond MiviaRover}
  \begin{columns}[T,onlytextwidth]
    \column{0.55\textwidth}
      \textbf{Contract-first across an SW/FW boundary} is a recurring
      pattern, not a MiviaRover oddity:
      \begin{itemize}\setlength\itemsep{2pt}
        \item \emph{Automotive} — every modern ECU networks via DBC/AUTOSAR
              IDLs.
        \item \emph{AGV / AMR} — VDA~5050 over MQTT plays the same role at
              the fleet level.
        \item \emph{Drones} — MAVLink dialects formalise ground/vehicle
              boundaries.
      \end{itemize}

      \vspace{0.4em}
      \textbf{Take-away as a method.}\;
      Define the \emph{wire format} before you write a line of code on
      either side. Then both teams can iterate independently, and your CI
      can lint compatibility from the file.

    \column{0.45\textwidth}
      \textbf{Concrete benefits we got out of it.}
      \begin{itemize}\setlength\itemsep{2pt}
        \item Re-flashing the MCU never breaks ROS code.
        \item Re-arching the ROS side never breaks the firmware.
        \item Simulating the bus is a one-liner
              (\texttt{cansend can0 0A\#…}).
        \item New telemetry = new \texttt{BO\_} block + new parser case;
              \emph{no protocol design call}.
      \end{itemize}
  \end{columns}
\end{frame}
\end{comment}
% ---------------------------------------------------------------------------- %
\begin{frame}{Take-aways from Atto~5}
  \begin{enumerate}\setlength\itemsep{2pt}
    \item \textbf{The loop is closed.} Encoder ticks come back as
          \texttt{/odom} through the same chain, walked in reverse.
    \item \textbf{The watchdog is the only fancy bit.} Two thresholds,
          two reactions: \emph{soft} fault zeroes velocities, \emph{hard}
          fault returns \texttt{ERROR} and engages \texttt{fault\_stop\_}.
    \item \textbf{Forward kinematics is free.} The same controller you
          configured downstream gives you odometry upstream — no extra
          code.
    \item \textbf{The DBC \emph{is} the boundary.} 5 messages, 1
          file, 1\,Mbit/s. The rest is implementation.
  \end{enumerate}
\end{frame}
